\paragraph{AQ-136-PD-RIGIDITY: finitely many local comparisons determine each divided-power operation.}
Let $C,D\in\mathcal A_k$, fix $r\ge1$ and $m\ge1$,
and use the operations of AQ-136-PD-OPERATIONS.
There is $W=W(C,D,r,m)$ such that
\[
           G_{C,m,w}^{2r}(a)=G_{D,m,w}^{2r}(a)
                    \qquad(w\ge W)                       \tag{136.PR1}
\]
for every finite simplicial set $X$ and
$a\in H^{2r}(X;k)$.
If the additive group of $k$ is finitely generated,
then (136.PR1) holds for \emph{every} simplicial set
$X$, with the same independence of the weight cutoff
from $X$. In particular the all-$X$ assertion holds
for the original integral coefficient ring $k=\mathbb Z$.
One may take $D=B$, the explicit standard integral
theory. This asserts equality with its operation,
without replacing that operation in torsion cohomology
by an unproved formula for an ordinary cup power.

\begin{proof}
For $m=1$ the operation is the identity, so assume
$m\ge2$ and write $d=m!$, $s=2r$, and $t=ms$.
For a rational prime $\ell$ let
\[
       k_{(\ell)}=k\otimes_{\mathbb Z}\mathbb Z_{(\ell)},
       \qquad C_{(\ell)}=C\otimes_k k_{(\ell)},
                                                               \tag{136.PR2}
\]
and define $D_{(\ell)}$ likewise. These are flat
localizations. Cycles, cohomology and normalized
simplicial homology commute with them, so the two
localized theories are admissible over $k_{(\ell)}$
with the original weight cutoffs. Divided powers
base change to the specified coefficient algebra
$\Gamma_{k_{(\ell)}}(T)$. If this coefficient ring
is zero, the assertions involving it are vacuous.

By the checked p-local common-source theorem, for
each of the finitely many primes $\ell\mid d$ there
is an admissible theory $R_\ell$ over $k_{(\ell)}$
and strict PD maps
\[
                    C_{(\ell)}\longleftarrow R_\ell
                                      \longrightarrow D_{(\ell)}.
                                                               \tag{136.PR3}
\]
In their common sufficiently large weight range,
naturality of the comparisons and of divided powers
therefore identifies the operations of these two
localized theories on all simplicial sets. There
is one such cutoff for each $\ell\mid d$, and we
take their finite maximum along with the cutoffs
for (136.PD1)--(136.PD2) for $C$ and $D$.

We specify coefficient naturality even on an infinite
simplicial set, without asserting that sections
commute with localization there. The actual map
$C\to C_{(\ell)}$ induces a map of the cycle
Eilenberg--Mac Lane spaces
\[
          Z^s_w(C)\longrightarrow Z^s_w(C_{(\ell)}).
\]
On their only nonzero homotopy groups it is exactly
$k\to k_{(\ell)}$, since normalization and homology
commute with the flat localization. The canonical
cohomology comparison is the coefficient-normalized
representing comparison proved in AQ-136-MULTIPLICATIVE.
Thus the induced map on cohomology for any $X$ is
the ordinary change-of-coefficients map. The same
argument applies in degree $t$ and weight $mw$.
Because the actual coefficient map commutes with
divided powers, it follows that
\[
 G_{C,m,w}^{s}(a)\big|_{k_{(\ell)}}
       =G_{C_{(\ell)},m,w}^{s}
                             (a\big|_{k_{(\ell)}}),
                                                               \tag{136.PR4}
\]
and likewise for $D$. Only these natural maps are
used here; none is claimed to be an isomorphism yet.

Set
\[
                  z=G_{C,m,w}^{s}(a)-G_{D,m,w}^{s}(a)
                               \in H^t(X;k).
\]
The factorial identity (136.PD2), whose right-hand
side is the same weighted cup power for both
theories, gives $dz=0$. Equations (136.PR3)--(136.PR4)
show that the image of $z$ in
$H^t(X;k_{(\ell)})$ is zero for every $\ell\mid d$.

Suppose first that $X$ is finite. Its normalized
integral chain complex is a bounded complex of
finite free modules. Applying $\operatorname{Hom}(-,k)$
and flat localization therefore gives the natural
isomorphism
\[
 H^t(X;k)\otimes_{\mathbb Z}\mathbb Z_{(\ell)}
                         \simeq H^t(X;k_{(\ell)}).
                                                               \tag{136.PR5}
\]
Consequently $z$ has zero localization at every
prime dividing $d$. An element killed by $d$ and
with this property is zero: if its nonzero order
were $e\mid d$, choose a prime $\ell\mid e$.
Vanishing in the $\ell$-localization would imply
$u z=0$ for an integer $u$ prime to $\ell$, whereas
the order $e$ would then divide $u$, a contradiction.
This proves (136.PR1) for all finite $X$, with one
cutoff for all of them.

Now suppose the additive group of $k$ is finitely
generated. Let $Y=K(k,s)$ and let
$\iota\in H^s(Y;k)$ be its fundamental class.
The integral homology groups of $Y$ are finitely
generated in every degree. An exact primary input
is Hatcher (2004), Lemma 1.10, p.~15, applied to
the Serre class of finitely generated abelian
groups. Its proof starts with cyclic groups and
uses the path fibration and Lemma 1.9; in particular
the statement applies to $s=2r\ge2$ here.

For any space whose integral homology in degrees
$t,t-1$ is finitely generated, (136.PR5) still
holds. To see this without localizing an infinite
cochain complex, use the natural universal
coefficient short exact sequence
\[
 0\to\operatorname{Ext}_{\mathbb Z}^1(H_{t-1}(Y;\mathbb Z),k)
  \to H^t(Y;k)
  \to\operatorname{Hom}_{\mathbb Z}(H_t(Y;\mathbb Z),k)\to0.
                                                               \tag{136.PR6}
\]
Each finitely generated homology group has a
resolution by two finite free abelian groups.
Both Hom and Ext in (136.PR6) therefore commute
with localization of $k$. Flatness localizes the
short exact sequence itself, and the natural
diagram with the coefficient-$k_{(\ell)}$ sequence
has isomorphisms at both ends. Its middle map is
an isomorphism by the short five lemma, proving
(136.PR5) for this particular universal space.

Apply the already established local comparison
(136.PR4) to $X=Y$ and $a=\iota$. Its difference
$z$ is $d$-torsion and is now detected by the
localizations (136.PR5). Thus $z=0$ on the
universal fundamental class. Every element of
$H^s(X;k)$ is the pullback of $\iota$ along a map
$X\to K(k,s)$, by representability. Naturality
then proves (136.PR1) for arbitrary $X$. No
additional cutoff depends on the representing
map or on $X$. This proves the all-$X$ assertion.
\end{proof}

\paragraph{AQ-136-LOCALIZATION-CAUTION: torsion cohomology on an arbitrary infinite space need not be detected by all coefficient localizations.}
The finite-space or universal finite-generation step
in the preceding proof is necessary to its reasoning.
Here is an explicit warning, unrelated to a counterexample
to original common-source existence.
Fix a prime $p$, let $A=\mathbb Z_{(p)}$, choose $s\ge2$,
and put $X=K(A,s)$. There is a nonzero class
\[
             z\in H^{s+1}(X;\mathbb Z),\qquad pz=0,
                                                               \tag{136.PR7}
\]
whose image in $H^{s+1}(X;\mathbb Z_{(\ell)})$ is
zero for \emph{every} rational prime $\ell$.

\begin{proof}
The first nonzero reduced integral homology of $X$
is $H_s(X;\mathbb Z)=A$. Universal coefficients
therefore give
\[
 H^s(X;\mathbb Z)=\operatorname{Hom}_{\mathbb Z}(A,\mathbb Z)=0,
 \qquad H^s(X;\mathbb Z/p)=\operatorname{Hom}_{\mathbb Z}
                                      (A,\mathbb Z/p)=\mathbb Z/p.
                                                               \tag{136.PR8}
\]
For the first equality, the image of $1$ under a
homomorphism $A\to\mathbb Z$ must be divisible by
every power of any fixed prime distinct from $p$,
and hence must be zero. Such a homomorphism is
then zero on every fraction. For the second, the
value at $1$ determines the map, and reduction
modulo $p$ extends to $A$ because its denominators
are prime to $p$.

The coefficient sequence
$0\to\mathbb Z\xrightarrow{p}\mathbb Z\to\mathbb Z/p\to0$
has a connecting map whose restriction from the
second group in (136.PR8) is injective, because
the first group is zero. Let $z$ be the image of
the nonzero reduction homomorphism. It is nonzero
and killed by $p$. After replacing integral
coefficients by $\mathbb Z_{(p)}$, that reduction
homomorphism lifts to the identity
$A\to\mathbb Z_{(p)}$ in $H^s$. Its connecting
class is consequently zero. Naturality of the
coefficient long exact sequence proves that the
image of $z$ is zero at $p$. For $\ell\ne p$,
multiplication by $p$ is an automorphism of the
entire coefficient cochain complex with values
in $\mathbb Z_{(\ell)}$ and hence of its cohomology.
The image of a $p$-torsion class there must be
zero. This proves (136.PR7) and the claimed
failure of naive detection.
\end{proof}

\paragraph{Scope of the rigidity result.}
For each fixed divided-power order and form degree,
the common-source results already proved at finitely
many coefficient primes force equality of the
cohomology operations on finite simplicial sets
over every $k$, and on all simplicial sets over
$\mathbb Z$ and more generally over a ring with
finitely generated additive group. This closes
that proposed operation-level distinction in
those scopes. It does not construct an integral
strict common source, a compatible tower of
finite-stage sources, or a contraction of all
ordinary-localization loops.

\paragraph{Exact additional reference.}
Hatcher, A. (2004, January 5). \emph{Spectral sequences
in algebraic topology}, Chapter 1: The Serre spectral
sequence. Author manuscript.
\url{https://pi.math.cornell.edu/~hatcher/SSAT/SSch1.pdf}.
Lemma 1.10 is stated on p.~15; its complete proof on
pp.~15--16 was directly inspected. Lemma 1.9, p.~15, supplies its fibration
step. Universal coefficients and representability are
the previously inspected Hatcher (2002), Theorems 3.2
and 4.57. Those primary results are cited inputs, not
new results of this module.
